\section{A chain model for contacts of arbitrary order}
\label{sec:chain-model-v169}
For $a\geq2$ consider the marked family over $S_a=\Spec k[b,c]$
\begin{equation}\label{eq:marked-family-v169}
 [F_0:F_1:F_2]=[t^a:b s^a:cst^{a-1}].
\end{equation}
Let $T_a$ be the closure of its Hilbert graph over $b\ne0$ in
$S_a\times\Hilb^{P_a}(X)$. This is a strict restriction of the
coefficient graph. The full fibre product with $\Gamma_a$ can have
additional vertical components and is not the definition of $T_a$.

\begin{theorem}[The all-order chain model]\label{thm:chain-model-v169}
For $m\geq m_a$ set
\[
 C=\frac{(a-1)m(m+1)}2,\quad
 E_j=m-j+1\ (1\leq j<a),\quad
 E_a=\frac{(m-a+1)(m-a+2)}2.
\]
The restriction of the evaluation ideal is exactly
\begin{equation}\label{eq:all-order-factor-v169}
 I_q(M_{a,m})\OO_{S_a}
   =b^C J_{a,m},\qquad
 J_{a,m}=\prod_{j=1}^a(b,c^j)^{E_j}.
\end{equation}
Its divisorial content is $b^C$. The ideals $J_{a,m}$ and all their
powers are integrally closed, and
\begin{equation}\label{eq:chain-blowup-v169}
 T_a=\operatorname{Bl}_{J_{a,m}}S_a
     =\operatorname{Bl}_{K_a}S_a,
 \qquad K_a=\prod_{j=1}^a(b,c^j).
\end{equation}
This surface is smooth. It is obtained by blowing up the origin and
then, $a-1$ times, the intersection of the newest exceptional curve
with the strict transform of $b=0$. Its exceptional divisors have
valuations
\[
 v_i(b)=i,\quad v_i(c)=1\qquad(1\leq i\leq a).
\]
The universal curve admits the explicit flat presentations below.
The theorem classifies the marked family \eqref{eq:marked-family-v169}
for every $a$; it does not identify the whole fibre of $\Gamma_a\to B_a$
with the fibre of this surface.
\end{theorem}

\subsection{A row factorization, without enumerating minors}
\begin{proof}[Proof of \eqref{eq:all-order-factor-v169}]
Every column has one nonzero entry. If $h,l$ are the exponents of the
second and third forms, that entry is $b^hc^l$ in row
$n=i+a(m-h-l)+(a-1)l$, $0\leq i\leq m$. A nonzero full minor chooses
one column in each row; columns belonging to distinct rows are
necessarily distinct. The determinant ideal is therefore the product
of the row ideals $Q_n$.

Put $d=am-n$. Rows with $n\geq am$ have $Q_n=(1)$. For $d\geq0$ set
\[
 h_0=\max\left(0,\left\lceil\frac{d-m}{a-1}\right\rceil\right),
 \qquad r=d-ah_0.
\]
For a fixed $h$, the smallest feasible $l$ is $\max(0,d-ah)$;
feasibility requires $h\geq h_0$. Removing dominated monomials gives
\begin{equation}\label{eq:row-factor-v169}
 Q_n=
 \begin{cases}
 (b^{h_0}),&r\leq0,\\
 b^{h_0}(b,c^a)^u(b,c^\rho),&r=au+\rho>0,\quad0\leq\rho<a,
 \end{cases}
\end{equation}
where the last factor is omitted for $\rho=0$. In fact the minimal
list in the second case is
$b^{h_0+v}c^{\max(0,r-av)}$, $0\leq v\leq\lceil r/a\rceil$.
Expanding the right side of \eqref{eq:row-factor-v169} and discarding
dominated terms gives precisely this list.

The sum of the $h_0$ is
$(a-1)(1+\cdots+m)=C$. Among all rows, a positive residual value
$r\in\{1,\ldots,m\}$ occurs
\[
 \nu(r)=\min(a,m-r+1)
\]
times. One occurrence comes from $h_0=0$; for $h_0=h\geq1$ it comes
from $m-a+2-h\leq r\leq m-h$. Counting the residual factors gives
\[
 \sum_{\substack{r\equiv j\pmod a\\1\leq r\leq m}}\nu(r)=m-j+1
 \quad(1\leq j<a),
 \qquad
 \sum_{r=1}^m\nu(r)\left\lfloor\frac ra\right\rfloor
       =\sum_{n=a}^m(n-a+1)=E_a.
\]
For the first equality write $n=r+v$, $0\leq v<a$, and count each
$n\in\{j,\ldots,m\}$ once. For the second, at a fixed $n$ sum
$\lfloor r/a\rfloor$ over
$\max(1,n-a+1)\leq r\leq n$; the answer is $\max(0,n-a+1)$.
This proves the ideal identity. The argument in fact works for every
$m\geq a$; the regularity bound is used only to identify its graph
with the Hilbert graph.
\end{proof}

\subsection{The Rees algebra and its rays}
\begin{proposition}[The normal Rees algebra]\label{prop:rees-v169}
Put $\eta_i=\sum_{j=1}^aE_j\min(i,j)$. Then
\begin{equation}\label{eq:rees-semigroup-v169}
 \mathcal R(J_{a,m})=
 k[b^u c^v T^n:\ u,v,n\geq0,\quad
                 iu+v\geq n\eta_i\ (1\leq i\leq a)].
\end{equation}
The right side is already normal, not merely the normalization of the
Rees algebra. For the raw ideal in \eqref{eq:all-order-factor-v169}
replace $u\geq0$ by $u\geq nC$ and the last inequalities by
$iu+v\geq n(iC+\eta_i)$. The exceptional Rees valuations of the
primitive ideal are exactly $v_1,\ldots,v_a$.
\end{proposition}
\begin{proof}
For a prescribed integer exponent $u$ of $b$, the least exponent of
$c$ in a product of the factors $(b,c^j)^{E_j}$ is obtained by replacing
$c^j$ by $b$ first for the largest $j$. Thus every integer abscissa
on the lower Newton boundary has an integer ordinate achieved by an
actual monomial generator. Its successive slopes are $-a,\ldots,-1$,
each with positive integral horizontal length $E_a,\ldots,E_1$.
The integral lattice points above this boundary are exactly the
monomials of the ideal. Replacing $E_j$ by $nE_j$ proves the same
assertion for every power. The half-space description of that boundary
is $iu+v\geq\eta_i$; adjoining the degree variable proves
\eqref{eq:rees-semigroup-v169}. A semigroup cut out by these integral
linear inequalities is saturated in its group. Its semigroup algebra
is normal. The compact facets have primitive inward normals $(i,1)$,
and their positive lengths show that none is redundant. They give
exactly the exceptional prime divisors of the normalized blow-up.
The shift by $b^{nC}$ proves the raw formula. The raw ideal also records
the nonexceptional strict transform of $b=0$; this is its content, not
an additional exceptional ray.
\end{proof}

\begin{proof}[Proof of the surface assertions in Theorem~\ref{thm:chain-model-v169}]
The graph Rees lemma and removal of $b^C$ give the first equality in
\eqref{eq:chain-blowup-v169}. The normal fan of the computed Newton
polygon has, in order, the rays
\[
 (1,0),(a,1),(a-1,1),\ldots,(1,1),(0,1).
\]
Every adjacent determinant is one. The same fan is obtained when all
$E_j$ are replaced by one, proving the second equality and smoothness.
Successive subdivisions insert $(1,1),(2,1),\ldots,(a,1)$ and give
the stated sequence of point blow-ups. Equivalently, a product of
positive powers of nonzero ideals is invertible on an integral test
scheme exactly when every factor is invertible as a fractional ideal;
the blow-up universal property gives the same identification.
\end{proof}

There are $a+1$ affine charts. The first and last are
\[
 T_0=\Spec k[b,k],\quad c=bk;
 \qquad T_a=\Spec k[c,z],\quad b=c^az.
\]
For $1\leq i<a$ the intermediate chart is
\begin{equation}\label{eq:chain-charts-v169}
 T_i=\Spec k[e,h],\qquad c=eh,\quad b=e^{i+1}h^i,
 \qquad e=b/c^i,\quad h=c^{i+1}/b.
\end{equation}
Adjacent transition formulas are
\begin{equation}\label{eq:chain-transitions-v169}
 \begin{gathered}
 k=e^{-1},\quad b=e^2h \quad(T_0,T_1),\\
 e_{i+1}=h_i^{-1},\quad h_{i+1}=e_i h_i^2
                  \quad(1\leq i<a-1),\\
 c=e_{a-1}h_{a-1},\quad z=h_{a-1}^{-1}
                  \quad(T_{a-1},T_a).
 \end{gathered}
\end{equation}
All further overlaps are compositions of these formulas.

\subsection{Low-degree equations for the universal curve}
In the source chart $s=1$, put $x=t/s$. Target coordinates $F,G,R$
remain homogeneous. On $T_0$ use
\begin{equation}\label{eq:curve-first-v169}
 R-kx^{a-1}G,\qquad x^aG-bF.
\end{equation}
On $T_i$, $1\leq i<a$, use
\begin{equation}\label{eq:curve-middle-v169}
 \begin{split}
 A&=xR-ehF,\\
 B_j&=x^{a-j}F^{j-1}G-e^{i+1-j}h^{i-j}R^j\quad(1\leq j\leq i),\\
 C_i&=R^{i+1}-h x^{a-i-1}F^iG.
 \end{split}
\end{equation}
On $T_a$ use
\begin{equation}\label{eq:curve-last-v169}
 A=xR-cF,\qquad
 B_j=x^{a-j}F^{j-1}G-zc^{a-j}R^j\quad(1\leq j\leq a).
\end{equation}
The ideal sheaves on $X$ are obtained by bihomogenizing these equations
and saturating by the source and target irrelevant ideals. In the
first chart a convenient global list adds $tR-bksF$ to the
bihomogenizations of \eqref{eq:curve-first-v169}; this prevents a
spurious component at $s=0$. In the other charts the listed
bihomogenizations already give the following two-source-chart cover.

\begin{proposition}[Flatness, exhaustion, and overlaps]\label{prop:curve-charts-v169}
Equations \eqref{eq:curve-first-v169}--\eqref{eq:curve-last-v169} and
\eqref{eq:chain-transitions-v169} define the flat universal embedded
curve over $T_a$. They cover all of that curve, including source
infinity. Each intermediate chart has a monic Gr\"obner basis over
$k[e,h]$, and the last chart has one over $k[c,z]$. Thus no
coefficient stratum is excluded by a division of a leading coefficient.
\end{proposition}
\begin{proof}
For the intermediate chart choose positive rational weights
\[
 \operatorname{wt}(F)=\operatorname{wt}(G)=1,\quad
 \operatorname{wt}(x)=\alpha,\quad
 \operatorname{wt}(R)=1+\rho\alpha,
 \qquad \frac{a-i-1}{i+1}<\rho<\frac{a-i}{i},\quad\alpha>0.
\]
Refine by any monomial order. The leading monomials of
\eqref{eq:curve-middle-v169} are respectively
\[
 xR,\qquad x^{a-j}F^{j-1}G\ (1\leq j\leq i),\qquad R^{i+1}.
\]
Buchberger reductions are uniform in the coefficient ring. The
$S$-polynomial of $A,B_j$ is $-ehB_{j+1}$ for $j<i$, and is $eC_i$
for $j=i$. The one for $A,C_i$ is $hFB_i$. For $j<l$, that of
$B_j,B_l$ is
\[
 e^{i+1-l}h^{i-l}R^j
       \big((xR)^{l-j}-(ehF)^{l-j}\big),
\]
and reduces by $A$. The remaining pairs have relatively prime leading
monomials. This proves the monic basis assertion and a free standard
monomial basis over $k[e,h]$.

For the last chart use lexicographic order $x>F>G>R$. The same
reductions replace $eh$ by $c$ and the coefficient of $R^j$ by
$zc^{a-j}$; the pair $A,B_a$ has relatively prime leading monomials.
The first chart eliminates $R$ and leaves the monic hypersurface
$x^aG-bF$. These bases make the target-graded coordinate rings flat
over their parameter rings. Localizing in a target coordinate and
taking the degree-zero part preserves flatness.

For completeness put $y=s/t$ on the other source chart. On every
parameter chart the bihomogeneous equations reduce to
\begin{equation}\label{eq:source-infinity-v169}
 G=b y^aF,\qquad R=c yF.
\end{equation}
For example $A$ gives the second equation, and $B_1$ gives the first;
the other equations then vanish. Here $F$ is a unit on the projective
target. This is the entire family near source infinity and is flat.
On $s,t\ne0$ the two displayed presentations agree by substitution.

Over the dense parameter torus all the equations are the graph of
\eqref{eq:marked-family-v169}. Their flat coordinate algebras are
torsion-free over the integral parameter ring, so each is the
schematic closure of that generic graph. Substitution of
\eqref{eq:chain-transitions-v169} therefore identifies the ideal
sheaves on each overlap, not merely their reduced supports: both are
the same closure in the fixed $X$. The substitutions of $b,c$ and the
identity on $[F:G:R]$ are explicit, so no target automorphism is hidden
in this argument. The glued family is projective, flat, and has generic
polynomial $P_a$; connectedness gives that polynomial on every fibre.
Its Hilbert map agrees with the defining map of $T_a$ on a dense open.
The graph identification proves that it is the universal curve.
\end{proof}

\begin{proposition}[Boundary geometry and contractions]\label{prop:boundary-v169}
The scheme fibre of $T_a\to S_a$ at the origin is the reduced chain
$E_1\cup\cdots\cup E_a$, with
\[
 E_i^2=-2\ (i<a),\qquad E_a^2=-1,\qquad
 E_iE_{i+1}=1,
 \qquad K_{T_a/S_a}=\sum_{i=1}^a iE_i.
\]
If $B',C'$ are the strict transforms of the two axes, then
\[
 \operatorname{div}(b)=B'+\sum iE_i,\qquad
 \operatorname{div}(c)=C'+\sum E_i.
\]
The Picard group is freely generated by the exceptional classes. A
relative divisor $D=\sum d_iE_i$ is nef precisely when
\[
 d_{i-1}-2d_i+d_{i+1}\geq0\ (1\leq i<a),\qquad
 d_{a-1}-d_a\geq0,\qquad d_0=0.
\]
Deleting rays gives the toric contraction morphisms. In particular,
contracting the intermediate chain between $(i,1)$ and $(j,1)$ gives
an $A_{j-i-1}$ singularity when $j-i>1$.
\end{proposition}
\begin{proof}
In an intermediate chart $(b,c)=(e^{i+1}h^i,eh)=(eh)$; at the end
charts its pullback is $(b)$ or $(c)$. Thus the fibre is reduced, with
ordinary crossings. The self-intersections follow either from the
successive point blow-ups or the relation between adjacent primitive
rays. The Jacobian in \eqref{eq:chain-charts-v169} is
$e^{i+1}h^i$, and the axis orders are the stated valuations. Each
point blow-up adds a free exceptional generator to the Picard group
of the preceding smooth surface, starting with $\Pic(\A^2)=0$.
The only contracted complete curves are the $E_i$, so intersecting
with them gives the nef inequalities. Coarsening a fan gives a proper
birational toric morphism. The two rays $(i,1),(j,1)$ have determinant
$j-i$ and lie at height one; the resulting affine surface is the
cyclic quotient, equivalently the hypersurface $uv=w^{j-i}$.
\end{proof}
The unweighted ideal $K_a$ is degree-independent and has exactly the
required Newton rays, but is not an inclusion-minimal graph ideal.
The powers $K_a^n$ already preclude such a minimality assertion. Its
boundary chain has a modular meaning for the marked family: it is the
universal modification on integral tests on which every $(b,c^j)$
becomes invertible. This is not a proposed intrinsic boundary complex
for all pencils of all contact types.
